\documentclass[12pt,a4paper]{exampaper}
\examname{高等数学分析 II}

\begin{document}

\examtitle{东南大学试卷（A 卷）}

\vspace{0.8em}

\begin{center}
  \examfield[3.4cm]{课程名称}{高等数学分析 II}\hspace{0.4em}%
  \examfield[2.2cm]{课程代码}{B07M1041}\hspace{0.4em}%
  \examfield[1.8cm]{考试学期}{2025-2026-2}\\[0.4em]
  \examfield[4.4cm]{适用专业}{工科类}\hspace{0.4em}%
  \examfield[2.0cm]{考试形式}{闭\hspace{0.5em}卷}\hspace{0.4em}%
  \examfield[1.9cm]{考试时长}{150\hspace{0.4em}分钟}
\end{center}

\vspace{0.8em}

\examscoretable[4]

\vspace{0.6em}

\examsection{填空题（本题共 9 小题，每小题 4 分，满分 36 分）}
\begin{examquestions}
  \examquestion 已知直线 $L_1: \frac{x-2}{1}=\frac{y+1}{m}=\frac{z-3}{2}$ 与直线 $L_2: x-1=\frac{y-2}{3}=\frac{z+1}{-1}$ 垂直，则 $m=$____
  \examquestion 曲面 $z=xy+\ln(1+x^2+y^2)$ 在点 $(0,1,0)$ 处的切平面方程为____
  \examquestion 设 $z=\arctan\frac{x+y}{1-xy}$（$xy\neq 1$），则 $dz|_{(0,0)}=$____
  \examquestion 设 $D=\{(x,y)\,|\,|x|+|y|\le 1\}$，则 $\iint_D |x+2y|d\sigma=$____
  \examquestion 设曲线 $L: x^2+y^2=4x$（$y\ge 0$），则 $L$ 的形心纵坐标 $\bar{y}=$____
  \examquestion 曲面 $z=x^2+y^2$ 位于柱面 $x^2+y^2=1$ 内的部分的面积为____
  \examquestion 若曲线积分 $\int_L \frac{(x+ay)dx+ydy}{x^2+y^2}$（$L$ 不绕过原点）与路径无关，则常数 $a=$____
  \examquestion 幂级数 $\sum_{n=1}^{\infty}\frac{n}{2^n}x^n$ 的收敛域为____
  \examquestion 设 $f(x)=|\sin x|$，$S(x)$ 是 $f(x)$ 的 Fourier 级数的和函数，则 $S\left(-\frac{\pi}{2}\right)=$____
\end{examquestions}

\vspace{0.6em}

\examsection{计算题（本题共 5 小题，每小题 7 分，满分 35 分）}
\begin{examquestions}[0pt]
  \examquestion 求均匀半球体 $\Omega: x^2+y^2+z^2\le R^2,\ z\ge 0$（体密度为 1）的质心坐标。
  \examanswerspace{7cm}
  \examquestion 求函数 $f(x,y)=x^2+2y^2-x^2y^2$ 在区域 $D=\{(x,y)\,|\,x^2+y^2\le 4,\ y\ge 0\}$ 上的最大值与最小值。
  \examanswerspace{7cm}
  \examquestion 判断级数 $\sum_{n=1}^{\infty}\frac{(-1)^n}{\sqrt{n}+\ln n}$ 的敛散性；若收敛，判断是绝对收敛还是条件收敛。
  \examanswerspace{7cm}
  \examquestion 计算曲线积分 $\int_L (2x-y+1)dx+(x+3y)dy$，其中 $L$ 为从 $O(0,0)$ 沿 $y=x^2$ 到 $A(1,1)$ 的一段弧。
  \examanswerspace{7cm}
  \examquestion 求幂级数 $\sum_{n=1}^{\infty}\frac{x^n}{n\cdot 2^n}$ 的收敛域与和函数。
  \examanswerspace{7cm}
\end{examquestions}

\vspace{0.6em}

\examsection{简答题（本题共 2 小题，每小题 8 分，满分 16 分）}
\begin{examquestions}[0pt]
  \examquestion 已知 $L$ 为闭曲线 $\begin{cases}x^2+y^2+z^2=4\\ x+y=1\end{cases}$，从 $z$ 轴正向往负向看去为逆时针方向，求曲线积分 $\oint_L (y+z)dx+(z+x)dy+(x+y)dz$。
  \examanswerspace{10cm}
  \examquestion 计算曲面积分 $\iint_\Sigma (x^2-y)dy\wedge dz+(y^2-z)dz\wedge dx+(z^2-x)dx\wedge dy$，其中 $\Sigma$ 为球面 $x^2+y^2+z^2=1$ 的外侧。
  \examanswerspace{10cm}
\end{examquestions}

\vspace{0.6em}

\examsection{证明题（本题共 2 小题，满分 13 分）}
\begin{examquestions}[0pt]
  \examquestion 证明：函数项级数 $\sum_{n=1}^{\infty}(-1)^n\frac{x^n+n^3}{3^n}$ 在区间 $[-1,1]$ 上一致收敛。
  \examanswerspace{8cm}
  \examquestion 设 $a_n>0$（$n=1,2,\cdots$），级数 $\sum_{n=1}^{\infty}a_n$ 收敛。\\ (1) 证明级数 $\sum_{n=1}^{\infty}\sqrt{a_na_{n+1}}$ 收敛；\\ (2) 举例说明级数 $\sum_{n=1}^{\infty}\frac{a_n}{n}$ 可能发散。
  \examanswerspace{8cm}
\end{examquestions}

\end{document}
